\renewcommand{\thechapterdisplay}{Vigésimo oitavo dia}
\renewcommand{\thechapterrunner}{Como evitar o Purgatório}
\chapter[Vigésimo oitavo dia --- Como evitar o Purgatório]{Como evitar o Purgatório}
\section*{Pensando frequentemente no Purgatório}

Considerai que o pensamento do Purgatório dirige naturalmente o nosso espírito ao pensamento da morte e do julgamento, e por isso pode nos inspirar salutares reflexões. \textit{Pensai em vosso fim último}, diz o Espirito Santo, \textit{e não pecareis.} A lembrança do Purgatório também orienta o espírito à penitência e à mortificação. Em vista dessa masmorra tão tenebrosa, de suas correntes tão pesadas, de suas chamas tão vivas e tão penetrantes, de suas angústias tão cruéis e tão longas, em vista de suas inumeráveis vítimas que emitem gritos tão dolorosos, a alma se examina e decide: \textit{Eu vou enfim resolver as contas abertas há tanto tempo com Deus; vou tirar proveito dos dias que Sua misericórdia me concede ainda para satisfazer à Sua justiça; vou quitar as dívidas que me são tão fáceis de pagar com um pouco de generosidade e de amor. Sim eu quero a todo custo evitar os tormentos do Purgatório. Eu posso, eu devo, eu o farei com a minha boa vontade e a graça de Deus.}

Oxalá sempre tivéssemos essa verdade diante dos olhos! É impossível que não nos tornássemos santos e grande santos. A lembrança habitual do Purgatório retirará de nossas vidas uma multidão de faltas leves, nos inspirará à prática das mais sublimes virtudes e, em nossa última hora, nossa alma, purificada pela penitência, enfeitada com méritos, voará em direção às moradas eternas, sem tocar nas chamas do Purgatório. Ó, morte feliz! O belo triunfo! Ó Purgatório, que sublimes lições nos ensinas!

\section*{Rezando frequentemente pelas almas do Purgatório}

Os padres e os doutores pensam que aqueles que se interessam verdadeiramente pelas almas do Purgatório escaparão dessas chamas ou não ficarão nelas por muito tempo. Porque, dizem eles, a marca mais infalível da predestinação é salvar almas, pois Deus prometeu fazer-nos o mesmo bem que fizermos aos outros. \textit{Bem-aventurados os misericordiosos, porque eles alcançarão misericórdia.} Além disso, não poderemos nós contar com a gratidão das almas que tivermos libertado? Poderiam elas mostrarem-se menos sensíveis e menos caridosas que nós? Na hora da nossa morte e do nosso julgamento, elas virão e estarão lá como protetoras, como testemunhas de defesa, para fazer pender a balança para o lado da misericórdia. Elas frustrarão as armadilhas do espírito infernal, acalmarão por suas constantes súplicas a cólera de nosso Juiz e nos obterão a mais preciosa das graças: uma santa morte. \textit{Eu não me lembro}, diz Santo Agostinho, \textit{de jamais haver lido sobre alguém que rezasse de bom grado pelos defuntos que teve uma morte ruim ou até mesmo duvidosa.}

Que meio quase seguro, alma cristã, de evitar os rigores do Purgatório! Sigamos, portanto, o conselho do Evangelho: façamos amigos, para que no momento da nossa morte, eles nos recebam nas moradas eternas.\footnote{(Lucas 16:9) \textit{Eu vos digo: fazei-vos amigos com a riqueza injusta, para que, no dia em que ela vos faltar, eles vos recebam nos tabernáculos eternos.}} Nossos irmão falecidos estão agora em necessidade, mas com um pouco da nossa ajuda, eles subirão ao Céu e nos abrirão em breve a porta para nós. Sim, libertemo-los do abismo do Purgatório e eles nos impedirão de cair ali. É relatado de Santa Margarida de Cortona, que em sua morte, todas as almas que ela havia libertado vieram recebê-la em triunfo com cantos e louvores.


\section*{Exemplo}

Conta-se que uma pessoa particularmente devota das almas do Purgatório e que havia dedicado sua vida a aliviá-las, tendo chegado à hora de sua morte, foi atacada com furor pelo demônio que a via a ponto de escapar-lhe. Parecia que todo o abismo conspirava contra ela, as hordas infernais cercavam-na. A moribunda lutava há algum tempo em meio aos mais penosos esforços, quando de repente viu entrar em seu aposento uma multidão de personagens desconhecidos, mas resplandecentes de beleza, que colocaram em fuga o demônio, e aproximando-se de seu leito, dirigiram-lhe exortações e consolações celestes. Exalando então um profundo suspiro e repleta de alegria, disse: \textit{Quem sois vós? Quem sois vós, por favor, vós que me fazeis tanto bem?} --- \textit{Nós somos os habitantes do Céu, que os vossos sufrágios conduziram à bem-aventurança, e viemos agora à nossa vez e por gratidão, vos ajudar a cruzar o limiar da eternidade, vos resgatar deste lugar de angústias, e vos introduzir nas alegrias da Cidade Santa.} Por essas palavras, um sorriso iluminou o rosto da moribunda, seu olhos se fecharam e ela adormeceu na paz do Senhor. Sua alma, branca e pura como uma pombinha, ao se apresentar ao soberano Juiz, encontrou tantos protetores e advogados quantas almas ela havia libertado e, reconhecida digna da glória, ela entrou no Paraíso como em triunfo, no meio de aplausos e de bênçãos de todos aqueles que ela havia tirado do Purgatório. Possamos um dia ter a mesma felicidade!


\vspace{\baselineskip}\noindent \textbf{Oremos.} Não permitais, ó meu Deus, que eu afaste de meu espírito, por uma errônea sensibilidade, o salutar pensamento do Purgatório. Gravai-o profundamente em meu coração, como um poderoso meio de preservar a mim mesmo do Purgatório e de me colocar em movimento, em auxílio das pobres almas que estão ali detidas. Como serei feliz se puder pôr um termo ao seu exílio e abrir-lhes a porta do Céu. Jesus, sede-lhes propício! Que descansem em paz!



\begin{center}
\textit{Requiescant in pace!}
\end{center}

Rezai agora uma dezena do terço e as orações no final deste livro.
